\subsection{Fases de desarrollo}
\label{section:fasdes}

\begin{enumerate}
  \item \textbf{Investigación y definición de conceptos} % Cap 2: Marco teórico
  \begin{itemize}
    \item Definición de conceptos básicos sobre el Juego de la Vida.
    \item Investigación de conceptos básicos sobre arquitecturas de cómputo paralelas.
    \item Definición de métricas de desempeño.
  \end{itemize}
  \item \textbf{Construcción y configuración del clúster} % Cap 3: Diseño e implementación
  \begin{itemize}
    \item Preparación de las computadoras.
    \item Instalación del sistema operativo a utilizar en las computadoras.
    \item Configuración de red, SSH, memoria compartida y asignación de tareas.
  \end{itemize}
  \item \textbf{Diseño e implementación del Juego de la Vida} % Cap 3: Diseño e implementación
  \begin{itemize}
    \item Diseño e implementación del Juego de la Vida en secuencial.
    \item Conversión paralela de la implementación secuencial del Juego de la Vida:
    \begin{itemize}
      \item Diseño de particionamiento del dominio.
      \item Diseño de las comunicaciones de los bordes de cada partición.
      \item Implementación paralela del Juego de la Vida.
    \end{itemize}
  \end{itemize}
  \item \textbf{Obtención y análisis de métricas} % Cap 4: Análisis de resultados
  \begin{itemize}
    \item Descripción de los escenarios de prueba.
    \item Comparación de los resultados del clúster contra una sola computadora.
  \end{itemize}
  \item \textbf{Redacción final} % Cap 5: Conclusiones y trabajo futuro
  \begin{itemize}
    \item Redacción de conclusiones del proyecto.
  \end{itemize}
\end{enumerate}